\documentclass[11pt]{article}
\usepackage[margin=1in]{geometry}
\usepackage{amsmath,amssymb,booktabs,hyperref}

\title{Exact Certification of the\\
11-Dimensional Kissing-Number Upper Bound}
\author{Ruturaj R Raval\\Independent Researcher\\
\texttt{ORCID: 0000-0003-4930-8981}}
\date{September 11, 2026}

\begin{document}
\maketitle

\begin{abstract}
We describe a proof-oriented implementation of the degree-17
Bachoc--Vallentin three-point semidefinite program for the
11-dimensional kissing number. The reported numerical objective is
$868.82650$. The certification target fixes the objective at
$86899/100$. A compact exact rational certificate at that objective passes
primary Julia verification and independent release-mode replay, establishing
$\tau_{11}\leq 868$. The package verifies all 70 matrix blocks
and every exact coefficient identity within the declared residual spaces.
The final prior-art refresh located no equivalent public certificate. The
release procedure independently replays the certificate before publishing
immutable tag-bound assets.
\end{abstract}

\section{Problem}

The kissing number $\tau_n$ is the largest number of nonoverlapping unit
spheres in $\mathbb{R}^n$ that can touch one central unit sphere. Equivalently,
it is the largest spherical code in $S^{n-1}$ whose distinct vectors have
inner product at most $1/2$.

The Bachoc--Vallentin three-point method gives semidefinite-programming upper
bounds for spherical codes. A degree-17 computation in dimension 11 reports
the numerical value $868.82650$, but the source paper states that its table
was not rigorously verified.

\section{Certification Target}

The project replaces objective minimization by the exact equality
\[
  \operatorname{objective}=\frac{86899}{100}=868.99.
\]
An exact feasible solution at this value implies
\[
  \tau_{11}\leq 868,
\]
because $\tau_{11}$ is an integer.

\section{Reduced Three-Point Bound}

Let $P_k^n$ denote the normalized Gegenbauer polynomial with
$P_k^n(1)=1$, and define
\[
 \Delta_s=\{(u,v,t)\in[-1,s]^3:
 1+2uvt-u^2-v^2-t^2\geq0\}.
\]
For nonnegative $a_0,\ldots,a_r$ and positive-semidefinite
$F_0,\ldots,F_d$, put
\[
 G(u)=\sum_{k=0}^r a_kP_k^n(u),\qquad
 H(u,v,t)=\sum_{k=0}^d
 \langle\bar Y_k^n(u,v,t),F_k\rangle.
\]
The reduced three-point bound states that
\[
 G(u)+3H(u,u,1)\leq-1\quad(u\in[-1,s])
\]
and
\[
 H(u,v,t)\leq0\quad((u,v,t)\in\Delta_s)
\]
imply
\[
 A(n,\arccos s)\leq
 M:=1+\sum_{k=0}^r a_k+\langle J,F_0\rangle.
\]

For completeness, let $C$ be a spherical code with $N=|C|$. Positivity of
the three-point kernels gives
\[
 T=\sum_{x,y,z\in C}
 H(x\mathbin{\cdot}z,y\mathbin{\cdot}z,x\mathbin{\cdot}y)\geq0.
\]
Separating triples by their number of distinct points and applying the two
displayed inequalities gives
\[
 T\leq NH(1,1,1)-N(N-1)
 -\sum_{k=0}^r a_k\sum_{x\ne y}P_k^n(x\mathbin{\cdot}y).
\]
Schoenberg positivity implies
\[
 \sum_{x\ne y}P_k^n(x\mathbin{\cdot}y)\geq-N.
\]
Since $H(1,1,1)=\langle J,F_0\rangle$, it follows that
$0\leq T\leq N(M-N)$ and hence $N\leq M$.

This is the reduced formulation of Dostert, de Laat, and Moustrou
\cite{DLM} and the mixed-degree formulation used by Leijenhorst and de Laat.
It is not described here as a literal transcription of Bachoc--Vallentin
Theorem 4.2, which also contains an auxiliary $2\times2$ block.

\section{Proof Architecture}

The certificate chain has four stages:
\begin{enumerate}
  \item reproduce the reported high-precision numerical bound;
  \item solve the fixed-objective feasibility problem with positive slack;
  \item project the numerical solution into the exact rational affine space;
  \item verify every affine identity and rational positive-semidefinite block
        with two independent implementations.
\end{enumerate}

The primary implementation uses Julia and Nemo together with the pinned
solver package \texttt{ClusteredLowRankSolver.jl}. The release certificate
uses a deterministic interchange format rather than Julia serialization.

The trivariate residuals lie in the symmetric polynomial space generated by
\[
  (u+v+t)^a(uv+ut+vt)^j(uvt)^k,
  \qquad a+2j+3k\leq 34.
\]
This space has dimension 1461. The canonical 1461-point evaluation matrix has
rank 1461 modulo 65521. Therefore it has full rank over the rationals, and
exact zero residuals at all canonical samples imply coefficient equality
within this declared space.

\section{Certificate Verification}

The rational problem generator and fixed-objective mode pass unit tests.
A degree-2 pilot in dimension 3 reproduces objective 14. A fixed-objective
pilot rounds to an exact rational feasible solution with zero affine
residuals and 16 positive-semidefinite blocks verified in exact arithmetic.
The exact verifier also rejects nonrational scalar types, wrong block layouts,
and objective mismatches, and it reconstructs the canonical problem rather
than trusting a serialized problem object.

A size audit rejects direct publication of the sampled affine system. The
1,663,458 nonzero three-point kernel terms alone require at least
93,153,648 bytes in that representation, and the complete expansion has an
upper bound of 505,535,295 terms. The replacement certificate stores each
exact rational symmetric upper triangle. The independent verifier proves
strict positivity of release matrices by directed fixed-point interval
Cholesky. A solver-produced degree-2 pilot reconstructs exactly in Julia and
passes complete independent Python coefficient replay.

The independent verifier uses only standard-library rational arithmetic. Its
56 tests cover the complete 70-block target layout, strict binary parsing,
exact objective and polynomial identities, malformed inputs, package-manifest
bindings, provenance recomputation, and resource limits. Nondegenerate
cross-language checks compare
Julia and Python on all 35 scalar kernels, three contractions for each of the
18 degree-17 \(F\) blocks, both interval SOS blocks, and three contractions
for each of the 15 three-point SOS families at degree 6.

The retained 256-bit dimension-11 run returned status \texttt{pdOpt} and
objective
\[
868.826499507902298742415116214534666\ldots.
\]
The accepted 384-bit fixed-objective recovery has objective error about
$1.17\times 10^{-68}$, maximum sampled affine residual about
$1.63\times 10^{-71}$, and minimum matrix eigenvalue about
$9.58\times 10^{-15}$. Portable exact affine projection and primary exact
verification now pass all 1,497 identities and all 70 strict-positive matrix
blocks. The resulting \texttt{KNWIT003} witness is 14,850,075 bytes and its
manifest is 12,458 bytes. Export-time captured-byte replay and a second
installed-file release-mode replay verify all 70 blocks, all 35 univariate
coefficients, and all 1,461 symmetric trivariate coefficients. Deterministic
compression produces a 14,565,324-byte package. The standalone replay took
907.16 seconds and 187,596,800 bytes maximum RSS, below the two-hour and
16 GiB limits. Together with the theorem bridge, these checks establish
\(\tau_{11}\leq868\).

\section{Release Procedure}

The mathematical, independent-verification, compactness, and resource gates
pass. The manuscript also builds reproducibly and passes visual inspection.
The release workflow performs a clean hosted replay, rebuilds the paper from
the tagged source, verifies every release-asset digest, and only then
publishes the immutable GitHub release. Zenodo archival follows that public
release and uses the same four byte-verified files.

\begin{thebibliography}{9}
\bibitem{BV}
C. Bachoc and F. Vallentin,
New upper bounds for kissing numbers from semidefinite programming,
\emph{Journal of the American Mathematical Society} 21 (2008), 909--924.

\bibitem{LL}
N. Leijenhorst and D. de Laat,
Solving clustered low-rank semidefinite programs arising from polynomial
optimization, arXiv:2202.12077.

\bibitem{CLL}
H. Cohn, D. de Laat, and N. Leijenhorst,
Optimality of spherical codes via exact semidefinite programming bounds,
arXiv:2403.16874.

\bibitem{DLM}
M. Dostert, D. de Laat, and P. Moustrou,
Exact semidefinite programming bounds for packing problems,
\emph{SIAM Journal on Optimization} 31 (2021), 1433--1458.

\bibitem{MO}
F. C. Machado and F. M. de Oliveira Filho,
Improving the semidefinite programming bound for the kissing number by
exploiting polynomial symmetry,
\emph{Experimental Mathematics} 27 (2018), 362--369.
\end{thebibliography}

\end{document}
